\documentclass[10pt,a4paper]{amsart}
\usepackage[margin=19mm]{geometry}
\usepackage{amsmath,amssymb,mathtools}
\usepackage[scr=rsfs,cal=euler]{mathalfa}
\usepackage{xfrac}
\usepackage[unicode,hidelinks,bookmarks=false]{hyperref}
\newcommand{\ie}{i.e.\ }
\newcommand{\cf}{cf.\ }
\newcommand{\resp}{respectively\ }
\newcommand{\Subsection}{\subsection}
\providecommand{\nobreakdash}{\nobreak\hbox{-}}
\setlength{\emergencystretch}{2em}
\multlinegap0pt
\usepackage{bm}
\usepackage{bbm}
\usepackage{dsfont}
\usepackage{xspace}
\usepackage{cancel}
\usepackage{mathtools}
\usepackage{amsthm}
\makeatletter
\@ifundefined{theorem}{\newtheorem{theorem}{Theorem}}{}
\@ifundefined{claim}{\newtheorem{claim}[theorem]{Claim}}{}
\@ifundefined{lemma}{\newtheorem{lemma}[theorem]{Lemma}}{}
\makeatother
\newcommand{\EE}{\mathbb{E}}
\newcommand{\PP}{\mathbb{P}}
\newcommand{\Sym}{\mathop{\mathrm{Sym}}}
\newcommand{\Unif}{\mathrm{Unif}}
\newcommand{\upperlinent}{\mathop{\overline{H}^{\mathrm{lin}}}}
\title[On the sampling entropy of permutons]{On the sampling entropy of permutons}
\author{Bal\'azs Maga}
\date{Source: arXiv:2503.18518v1 (CC BY 4.0)}
\begin{document}
\maketitle

\noindent\emph{Source and licence.} On the sampling entropy of permutons. Version arXiv:2503.18518v1 (CC BY 4.0), distributed under the Creative Commons Attribution 4.0 International licence, as recorded in the arXiv licence metadata for this exact version and shown on the abstract page https://arxiv.org/abs/2503.18518v1. Full rights notice: licenses/arxiv-2503.18518-CC-BY-4.0.txt.\par\smallskip

\noindent\textit{Editorial scope.} Counted unit: the Lemma whose source label is \texttt{lemma:random\_aut\_conv\_means} in the article (source0.tex lines 832--834, source label \texttt{lemma:random\_aut\_conv\_means}), delivered together with the article's own complete original proof of it (source0.tex lines 836--979, one \texttt{proof} environment of 144 lines). No mathematics is written, repaired or re-derived: statement and proof are the article's own text line for line. Required context retained uncounted: thm:random\_automorphism (lines 133--140); lemma:automorphism\_finite\_lin\_entropy (lines 648--657); eq:entropy\_decomp (lines 740--742); eq:rearranged (lines 767--769); eq:recursion\_alpha (lines 779--781); eq:disintegration (lines 825--827). Every definition, display and label of those lines is retained unchanged. Omitted from this package and not redistributed: 1--132, 141--647, 658--739, 743--766, 770--778, 782--824, 828--831, 835--835, 980--1107. Nothing inside a retained excerpt is omitted or altered: every retained body is delivered exactly as the article's lines are written, so no omission receipt of any kind accompanies this package. Citations: the retained excerpts carry no citation command at all, so no bibliography entry of the article is transcribed and no reference is redistributed. Reference adaptation: none was needed and none was made; every reference inside the delivered unit resolves inside it, so no unresolved reference or citation is produced. Printed slips kept literal: no printed word, symbol, hypothesis or number of the counted statement or of its proof was altered. Layout and class: the article's own class is \texttt{article} with packages including fullpage, graphicx, amsfonts, amsthm, amsmath, amssymb, mathabx, mathrsfs, mathptmx, comment, xcolor, biblatex, titling, hyperref; the wrapper is the lane's pinned \texttt{amsart} editorial wrapper at 10pt on A4, which restates the theorem-like environments the retained text needs and adds the article's own macro definitions that the retained text needs (\texttt{\textbackslash EE}, \texttt{\textbackslash PP}, \texttt{\textbackslash Sym}, \texttt{\textbackslash Unif}, \texttt{\textbackslash upperlinent}), with extra wrapper packages (bm, bbm, dsfont, xspace, cancel, mathtools). The delivered numbering is the unit's own. Assets: no retained span contains a figure environment, a graphics-inclusion command, a PDF-page inclusion, a table or a TikZ picture; no image file is copied into this package, the package declares no asset dependency and no image file is redistributed with it. Licence: version arXiv:2503.18518v1 of this article is distributed under the Creative Commons Attribution 4.0 International licence, as recorded in the arXiv licence metadata for this exact version and shown on the abstract page https://arxiv.org/abs/2503.18518v1. A full rights notice is delivered at licenses/arxiv-2503.18518-CC-BY-4.0.txt. Characters: every retained excerpt is pure ASCII, so the delivered engine, XeLaTeX with its default Latin Modern text font, reports no missing character.
\begin{theorem}\label{thm:random_automorphism}
    For the random permuton $\mu=\mu_F$ induced by $F\in \PP(\Sym(d))$, $H_n(\mu)/n$ is asymptotically log-periodic in base $d$ with limit
    \begin{equation} \label{eq:aut_thm_statement}
    A(x)=\sum_{r=-\infty}^{\infty}\left(\sum_{l=1}^{\infty}\rho_{l+1}\frac{1-d^{-l}}{(l+1)! \log d}\Gamma\left(l-\frac{2\pi i r}{\log d}\right)\right)\exp (2\pi i r x),
    \end{equation}
    for some nonnegative sequence $\rho_l=O(\log l)$ determined by $F$. The average $\int_{0}^{1}A(x)dx$ vanishes if and only if
    $F$ is a Dirac measure on $(1,2\dots ,d)$ or $(d,d-1,\dots ,1)$. 
\end{theorem}

\begin{lemma} \label{lemma:automorphism_finite_lin_entropy}
The followings hold for $X=X_\Pi$ and $\mu=\mu_f$:
\begin{enumerate}
    \item $X$ is a substitution-closed permutation class and $\pi_n(f)\in X$.
    \item If $\sigma\notin X$, then $t(\sigma, \mu_f)=0$ for any outcome.
    \item If $\sigma \in X$, then $t(\sigma, \mu_f)>0$ with probability 1.
    \item $X$ is a pattern-avoiding class.
    \item For some constant $K\geq 0$ depending only on $d$, $H_n(\mu_f)\leq Kn$, and hence $\upperlinent(\mu_f)<\infty$.
\end{enumerate}
\end{lemma}

    \begin{equation}\label{eq:entropy_decomp}
    X_n = Z_{n, m} + \sum_{i=1}^{d^m} \sum_{k_i=0}^{n}\binom{n}{k_i} d^{-mk_i}(1-d^{-m})^{n-k_i}H_{k_i}(\nu_i)=Z_{n, m} + \sum_{i=1}^{d^m}V_{i, n, m}.
    \end{equation}

    \begin{equation}\label{eq:rearranged}
    y_n = \rho_n + \frac{d^n}{d^{n-1}-1}\sum_{k=0}^{n-1}\binom{n}{k}d^{-k}(1-d^{-1})^{n-k} y_k=\rho_n + \sum_{k=0}^{n-1}\frac{\binom{n}{k}}{d^{n-1}-1}(d-1)^{n-k}y_k,
    \end{equation}

    \begin{equation}\label{eq:recursion_alpha}
        \alpha_l(n)=\sum_{k=0}^{n-1}\frac{\binom{n}{k}}{d^{n-1}-1}(d-1)^{n-k}\alpha_l(k),
    \end{equation}

\begin{equation} \label{eq:disintegration}
\PP_{\mu}(A)=\EE_{\tau|_m}(\PP_{\mu^{t|_m}}(A)),
\end{equation}

\begin{lemma}\label{lemma:random_aut_conv_means}
    For $\mu=\mu_{F,\Unif}$, $\EE H_n(\mu)/n$ converges.
\end{lemma}

\begin{proof}
    Recall that on the first level, $(\tau_i)_{i=1}^{d}$ are the random side lengths of the squares in the first approximation of $\mu$, being gap sizes in order statistics obtained from $U_1, \dots, U_{d-1}\sim \Unif(0, 1)$. Denote the resulting intervals of the partition by $I_1, \dots, I_d$. Fixing a first-level choice $(\tau_i)_{i=1}^{d} = (t_i)_{i=1}^{d}$, we can write in the spirit of \eqref{eq:entropy_decomp}
    \begin{equation}\label{eq:entropy_decomp2}
    H_n(\mu^{t|_1}) = Z_{n, 1}(\mu^{t|_1}) + \sum_{i=1}^{d} \sum_{k_i=0}^{n}\binom{n}{k_i} t_i^{k_i}(1-t_i)^{n-k_i}H_{k_i}(\nu_i).
    \end{equation}
    Moreover, $0\leq Z_{n, 1}(\mu) \leq 2d\log n$ for $n$ large.
    Note that here
    $$\binom{n}{k_i} t_i^{k_i}(1-t_i)^{n-k_i}=\PP(\text{upon sampling $n$ points from $\mu$, there are $k_i$ points with $x$-coordinate in $I_i$}).$$
    Taking expectation with respect to the random node labelling and nonfixed edge weights, we find due to $\nu_i$ and $\mu$ being equidistributed that
    \begin{equation}\label{eq:entropy_decomp2_expectation}
    \EE[H_n(\mu^{t|_1})]=\EE [Z_{n, 1}(\mu^{t|_1})] + \sum_{i=1}^{d}\sum_{k_i=0}^{n}\PP(\text{there are $k_i$ points in $I_i$})\cdot\EE[H_{k_i}(\mu)]
    \end{equation}
    Taking expectation with respect to the choice of $(t_i)$, due to the gap sizes being equidistributed and having the natural disintegration described in \eqref{eq:disintegration},
    $$\EE [H_n(\mu)] = \EE [Z_{n, 1}(\mu)] + d\sum_{k=0}^{n}\PP(\text{there are $k$ points in $[0, U_{(1)}]$})\cdot\EE [H_k(\mu)].$$
    Note that the distribution on the length $d$ order-dependent partitions of $n$ induced by the partition given by $U_{(0)}, \dots, U_{(d)}$ is uniform, thus we can easily express 
    \begin{equation}\label{eq:point_distribution}
    \PP(\text{there are $k$ points in $[0, U_{(1)}]$})=\frac{\binom{n-k+d-2}{d-2}}{\binom{n+d-1}{d-1}}
    \end{equation}
    by enumerating the corresponding combinations with repetition. 
    
    We check $d=2$ separately, as this case is deceptively simple: \eqref{eq:point_distribution} is just a uniform distribution on $\{0, 1, \dots, n\}$, hence
    \begin{equation}\label{eq:recursion_d=2}
    \EE [H_n(\mu)] = \EE [Z_{n, 1}(\mu)] + \frac{2}{n+1}\sum_{k=0}^{n}\EE [H_k(\mu)],
    \end{equation}
    or after rearrangement and writing $y_n=\EE [H_n(\mu)]$,
    $$y_n = \rho_n + \frac{2}{n-1}\sum_{k=0}^{n}y_k,$$
    where $0\leq \rho_n = \frac{n+1}{n-1}\EE [Z_{n, 1}(\mu)]\leq 4d\log n$. (Note that we have the same upper bound for any $d$ if $n$ is large enough.)
    
    Now in the same fashion as in proceeding from \eqref{eq:rearranged} to \eqref{eq:recursion_alpha} in the proof of Theorem \ref{thm:random_automorphism}, this implies
    $$y_n = \sum_{l=2}^{n}\alpha_l(n)\rho_l,$$
    where $\alpha_l(n)=0$ for $n<l$, $\alpha_l(l)=1$, and for $n>l$,
    $$\alpha_l(n)=\frac{2}{n-1}\sum_{k=0}^{n-1}\alpha_l(k).$$
    Applying this recurrence,
    $$\alpha_l(n+1)= \frac{2}{n}\alpha_l(n)+\frac{2}{n}\sum_{k=0}^{n-1}\alpha_l(k)=\frac{n-1}{n}\alpha_l(n)+\frac{2}{n}\alpha_l(n)=\frac{n+1}{n}\alpha_l(n),$$
    thus $\alpha_l(n)/n$ is constant for $n=l, l+1, \dots$. By plugging in $n=l$, we see that its value is $\frac{1}{l}$, yielding
    $$y_n/n = \sum_{l=2}^{n}\rho_l/l.$$
    It implies that $y_n/n$ is a nondecreasing sequence, whose limit is finite due to 5. in Lemma \ref{lemma:automorphism_finite_lin_entropy}. This concludes the proof for $d=2$.

    The rest of the proof is devoted to $d>2$. By \eqref{eq:point_distribution}, we find that $\alpha_l(n)$ satisfies the following recursion for $n>l$:
    \begin{equation}\label{eq:recursion_alpha_d>2}
    \alpha_l(n)= d\sum_{k=0}^{n}\frac{\binom{n-k+d-2}{d-2}}{\binom{n+d-1}{d-1}}\alpha_l(k),
    \end{equation}
    or after rearrangement,
    \begin{equation}\label{eq:recursion2_alpha_d>2}
    (n+d-1)_{d-1}\alpha_l(n)= d(d-1)\sum_{k=0}^{n}(n-k+d-2)_{d-2}\alpha_l(k),
    \end{equation}    
    where $(x)_m=\prod_{i=0}^{m-1}(x-i)$ is the falling factorial. Recalling
    $$y_n = \sum_{l=2}^{n}\alpha_l(n)\rho_l,$$
    proving the convergence of each $\alpha_l(n)/n$ separately is a promising direction. We note that we can quickly observe that $\alpha_l(n)=n$ (and hence, by linearity, $\alpha_l(n)=Cn$ for any $C$) is a solution. Notably, then the identity in question is
    $$(n+d-1)_{d}= d(d-1)\sum_{k=0}^{n}(n-k+d-2)_{d-2}\cdot k.$$     
    Now the left hand side is the number of ways of choosing an ordered list of length $d$ consisting of distinct elements drawn from a collection of size $n+d-1$. The right hand side counts the same quantity, decomposed based on the location of the second entry: if the second entry is at place $k+1$, there are $(n-k+d-2)_{d-2}\cdot k$ to choose the rest. Thus $\alpha_l(n)=Cn$ is a solution indeed, which is promising for our goal that $\alpha_l(n)$ grows linearly. We need to check essentially that all solutions independent from $Cn$ are sublinear.

    Before dealing with this recurrence relation or difference equation, it is enlightening to study the naturally corresponding differential equation, where the set of solutions can be understood in simpler terms. Replacing the sum by an integral and using the approximation $n+o(n)\approx n$, we can rewrite \eqref{eq:recursion2_alpha_d>2} as
    $$x^{d-1}f(x)=d(d-1)\int_{0}^{x}(x-\xi)^{d-2}f(\xi)d\xi.$$
    Differentiating $d-1$ times with respect to $x$, we obtain a Cauchy--Euler equation of a very simple form:
    \begin{equation}\label{eq:cauchy_euler}
    \left(x^{d-1}f(x)\right)^{(d-1)}=d! \cdot f(x).
    \end{equation}
    This is a linear differential equation of order $d-1$, so there are $d-1$ linearly independent solutions. One should look for solutions of the form $f(x)=x^r$, which is a solution if and only if
    $$\prod_{i=1}^{d-1}(r+i)=d!,$$
    or after writing $s=r+d-1$,
    $$(s)_{d-1}=d!$$
    Finding $d-1$ distinct roots for this polynomial equation gives a basis of the space of solutions of \eqref{eq:cauchy_euler}. If the solutions are not distinct, some extra care is needed, but we do not face this inconvenience. The following claim will have importance in the exact treatment of the recurrence relation as well:

    \begin{claim}\label{claim:falling_factorial_roots}
    The polynomial $p(s)=(s)_{d-1}-d!$ have $d-1$ distinct roots, one of them is $d$, and all other roots have real part smaller than $d$.
    \end{claim}

    \begin{proof}[Proof of Claim \ref{claim:falling_factorial_roots}]
        If $\Re s>d$ or $\Re s = d$ and $s\neq d$, then comparing the absolute values of the terms one-by-one in $(s)_{d-1}$ and $d!=(d)_{d-1}$ shows that $|(s)_{d-1}|>|d!|$. To check that the roots are distinct, observe that if $s\in (-1, d)$, then by the same one-by-one comparison of absolute values $|(s)_{d-1}|<|d!|$. Thus any real root of $p(s)$ is outside of the interval $(-1, d)$. On the other hand, the $d-1$ roots of $q(s)=(s)_{d-1}$ are $s=0, \dots, d-2$, and then by the Gauss--Lucas theorem, all roots of $q'=p'$ are in the interval $[0, d-2]$. Thus there are no common roots of $p, p'$, yielding the existence of $d-1$ distinct roots.
    \end{proof}

    By the previous claim, $\prod_{i=1}^{d-1}(r+i)=d!$ is solved by $r=1$ and all other roots have real part smaller than 1, i.e., any solution of \eqref{eq:cauchy_euler} is of the form $f(x)=Cx+o(1)$. This is the type of statement we want to achieve for the solution of the original recurrence.

    Having made these preliminary observations, we highlight the claim we have already formulated and which we will prove now rigorously:
    
    \begin{claim} \label{claim:alpha_l_n_conv}
    $\alpha_l(n)/n$ is convergent for any $l$. 
    \end{claim}

    \begin{proof}[Proof of Claim \ref{claim:alpha_l_n_conv}]
    Motivated by the way we solved the continuous version, recall the recurrence relation \eqref{eq:recursion_alpha_d>2}
    \begin{equation}\label{eq:recursion2_alpha_d>2_again}
    (n+d-1)_{d-1}\alpha_l(n)= d(d-1)\sum_{k=0}^{n}(n-k+d-2)_{d-2}\alpha_l(k),
    \end{equation}
    and apply discrete differentiation $d-1$ times to get rid of the "integral". Discrete differentiation of a sequence $a(n)$ just accounts for taking $\Delta a(n)=a(n+1)-a(n)$, or after multiple steps, by induction,
    $$\Delta^{(m)}a(n)= \sum_{i=0}^{m}(-1)^{m-i} \binom{m}{i}a(n+i).$$
    Applying this with $m=d-1$ to the right hand side of \eqref{eq:recursion2_alpha_d>2_again}, by definition
    \begin{equation}\label{eq:discrete_diff}
    \Delta^{(d-1)}\left(d(d-1)\sum_{k=0}^{n}(n-k+d-2)_{d-2}\alpha_l(k)\right) =d(d-1)\sum_{i=0}^{d-1}(-1)^{d-1-i} \binom{d-1}{i}\sum_{k=0}^{n+i}(n+i-k+d-2)_{d-2}\alpha_l(k).
    \end{equation}
    Observe that for any $i$, the second sum can be extended to run until $n+d-2$, as plugging in $k=n+i+1, \dots, n+d-2$ vanishes in the falling factorial. Making this extension, decomposing the first sum to $i=d-1$ and $i$ running from $0$ to $d-2$, and exchanging the sums, we see that \eqref{eq:discrete_diff} further equals
    $$d(d-1)\left((d-2)_{d-2}\cdot \alpha_l(n+d-1)+\sum_{k=0}^{n+d-2}\alpha_l(k)\sum_{i=0}^{d-1}(-1)^{d-1-i}\binom{d-1}{i}(n+i-k+d-2)_{d-2}\right).$$
    Observe that in this sum, the coefficient of $\alpha_l(k)$ is just the $(d-1)$th discrete derivative of the degree $d-2$ polynomial $(n-k+d-2)_{d-2}$, up to sign, hence it vanishes. Thus what remains is that the $(d-1)$th discrete derivative of the right hand side of \eqref{eq:recursion2_alpha_d>2_again} is just $d!\cdot\alpha_l(n+d-1)$. Simply plugging the definition of $\Delta$ into the left hand side, we see that taking these derivatives yields the following recurrence of order $d-1$:
    \begin{equation}\label{eq:final_recurrence}
    \sum_{i=0}^{d-1}(-1)^{d-1-i} \binom{d-1}{i}(n+i+d-1)_{d-1}\alpha_l(n+i)=d!\cdot\alpha_l(n+d-1).
    \end{equation}
    Finding $d-1$ linearly independent solutions solves this recurrence. The solutions naturally associated with the polynomial solutions of the corresponding differential equation are the hypergeometric solutions, i.e., we look for solutions of the form $\alpha_l(N+1)/\alpha_l(N)=S(N)$, where $S(N)$ is some rational function. We will be even more restrictive with our choice motivated by what we found in the continuous case, we will look for $S(n)=\frac{N+1}{N+x}$. 
    Then we have 
    $$\alpha_l(n+i)=\alpha_l(n)\cdot \frac{(n+i)_i}{(n+x+i-1)_i}.$$
    Plugging this into \eqref{eq:final_recurrence}, and multiplying by $(n+x+d-2)_{d-1}$, we can simplify by the common factors $(n+d-1)_{d-1}\cdot \alpha_l(n)$ to get a polynomial equation in $n$ and $x$, having degree $d-1$ in both variables:
    \begin{equation}\label{eq:recurrence_x_form}
    \sum_{i=0}^{d-1}(-1)^{d-1-i} \binom{d-1}{i}(n+i+d-1)_{i}\cdot (n+x+d-2)_{d-1-i}=d!.
    \end{equation}    
    Denote the polynomial on the left hand side by $P(x, n)$. The vital observation is that $P(x, n)$ is actually independent of $n$. As we already now that $\alpha_l(n)=n$ solves the recurrence, this $n$-independence holds for $x=0$. Moreover, note that upon plugging in $x=1,2 , \dots, d-1$ into the left hand side, we get
    $$\prod_{j=0}^{x-1}(n+d+x)\sum_{i=0}^{d-1}(-1)^{d-1-i}\binom{d-1}{i}(n+i+d-1)_{d-x}=\prod_{j=0}^{x-1}(n+d+x) \Delta^{d-1}(n+i+d-1)_{d-x}.$$
    This is the product of a degree $x-1$ polynomial and the $(d-1)$th derivative of a degree $d-x$ polynomial, which is zero unless $x=1$, when it is still constant. Thus $P(x, n)$ is independent of $n$ on each of the lines $x=0, 1, \dots, d$. Thus the coefficient of any positive power of $n$ vanishes for $x=0, 1, \dots, d$. But these coefficients are degree $d-1$ polynomials of $x$, hence it is possible only if they are identically zero. Thus $P(x, n)$ is independent of $n$ indeed, and we have already noted that $x=2, \dots, d-1$ are its roots. Recalling \eqref{eq:recurrence_x_form}, we can see that the leading coefficient is $(-1)^{d-1}$, thus 
    \begin{equation}\label{eq:x}
    P(x, n)=(-1)^{d-1}(x-2)_{d-1}=d!,
    \end{equation}
    or plugging in $s=-(x-d)$,
    $$(s)_{d-1}=d!.$$
    We understood this equation in Claim \ref{claim:falling_factorial_roots}: it has $d-1$ distinct roots, including $d$, and all other roots have real part smaller than $d$. As $d-s=x$, this yields that \eqref{eq:x} has $d-1$ distinct roots, including 0, and all other roots have real part larger than 0. Denoting these roots by $x_1, \dots, x_{d-1}$, we have obtained that
    $$\alpha_l^{(i)}(n) = \frac{(n)_{n-l}}{(n+x_i-1)_{n-l}}\alpha_l(l)=\frac{(n)_{n-l}}{(n+x_i-1)_{n-l}}$$
    is a solution for any $i$, $n\geq l$. It is a standard exercise to check that these $d-1$ solutions are linearly independent, thus they span the space of all solutions, and apart from $x_1=1$, all these are sublinear. Thus all solutions are of the form
    $$\alpha_l(n)=\sum_{i=1}^{d-1}c_i\alpha_l^{(i)}(n) =c_1n+o(n).$$
    Thus $\alpha_l(n)/n$ is convergent for any $l$.
    
    These roots yield the $d-1$ linearly independent solutions of \eqref{eq:recursion2_alpha_d>2_again}. 
    \end{proof}

    Now we can conclude the proof of Lemma \ref{lemma:random_aut_conv_means}. Recall that we need that
    \begin{equation}\label{eq:y_n_from_rho}
    y_n/n =\sum_{l=2}^{n}\frac{\alpha_l(n)}{n}\rho_l
    \end{equation}
    is convergent as $n\to\infty$. Due to the convergence of the coefficients $\alpha_l(n)/n$, it suffices to prove that for any $\varepsilon>0$ we can find $l_\varepsilon$ such that for every $n$
    $$\sum_{l=l_\varepsilon}^{n}\frac{\alpha_l(n)}{n}\rho_l<\varepsilon.$$
    To prove this, we observe that $\alpha_l(n)$ has a similar stochastic meaning as in the proof of Theorem \ref{thm:random_automorphism}. In this case, we do the following random process: starting from $n$ stones labelled by $[n]$, we drop $d-1$ separating bars independently uniform at random (they might end up in the same locations and they might fall before or after all the stones). Then in the resulting intervals, we repeat the same process independently until every stone is the sole inhabitant of its interval. We call this the {\it simple partition decay}. Then $\alpha_l(n)$ is the expected number of length $l$ intervals. Here
    \begin{equation}\label{eq:alpha_stochastic_meaning}
    \EE(\text{no. of $l$-long intervals})=\frac{1}{l}\sum_{i=1}^n\PP(\text{$i$ appears in an $l$-long interval}).
    \end{equation}
    Note that in contrast to the other proof, this quantity does not equal
    $\frac{n}{l}\PP(\text{1 appears in an $l$-long interval})$, as the probability $\PP(\text{$i$ appears in an $l$-long interval})$ depends on $i$. (This can be easily seen via the example that the event "1 appears in an $l$-long interval" is contained by the event "2 appears in an $l$-long interval if $l>1$".) Thus we need that for large enough $l_\varepsilon$
    $$\sum_{l=l_\varepsilon}^{n}\frac{1}{nl}\sum_{i=1}^n\PP(\text{$i$ appears in an $l$-long interval})\rho_l<\varepsilon.$$
    We partition the interval $[l_\varepsilon, n]$ into intervals $I_0, I_1, \dots, I_m$, where $I_j =[r^jl_\varepsilon, r^{j+1}l_\varepsilon]$ for some $r>1$. Summing over $l\in I_j$, by the bound on $\rho_l$, we get
    \begin{equation}\label{eq:prob_bound_partition_decay}
    \sum_{l\in I_j}\frac{1}{nl}\sum_{i=1}^n\PP(\text{$i$ appears in an $l$-long interval})\rho_l\leq \frac{4d \log \max I_j}{n \min I_j}\sum_{i=1}^{n}\sum_{l\in I_j}\PP(\text{$i$ appears in an $l$-long interval}).
    \end{equation}
    This sum of probabilities is the expected number of intervals encountered in the simple partition decay containing $i$ and having length in $I_j$, thus we must understand the evolution of this $i$-container length. We can observe that with probability exceeding 1/2, the length of the $i$-container is multiplied by at most $3/4$. (If $d=2$, and $i$ is in the middle of its container, this decay has precisely probability 1/2, and it is plain to see that otherwise it is even larger.) Thus if $r<4/3$, the number of $i$-containers with length in $I_j$ is stochastically dominated by $1+G$ for $G\sim \mathrm{Geo}(1/2)$. Hence the expected number of intervals encountered in the simple partition decay containing $i$ and having length in $I_j$ is at most $\EE(1+G)=3$. Hence \eqref{eq:prob_bound_partition_decay} can be bounded further by
    $$\frac{12d \log \max I_j}{\min I_j}$$
    for any $n$, and hence 
    $$\sum_{l=l_\varepsilon}^{n}\frac{1}{nl}\sum_{i=1}^n\PP(\text{$i$ appears in an $l$-long interval})\rho_l\leq \sum_{j=0}^{m}\frac{12d \log \max I_j}{\min I_j}=\frac{12d}{l_\varepsilon}\sum_{j=0}^{m}\frac{j+1}{r^j}.$$
    The convergence of this series assures that this quantity is indeed at most $\varepsilon$ for $l_{\varepsilon}$ large enough. This concludes the proof.
    \end{proof}

\end{document}
